% Last proposal background paragraph retained; theory is separated for the 2026 template.
\subsection{Theory and Related Work}
Closer studies delimit our question. Yang et al. plan cooperative USV seabed-survey coverage paths \cite{yang2023cooperative}. Bramblett and Bezzo coordinate coverage and task allocation using beliefs about disconnected robots \cite{bramblett2023epistemic}. Gordon et al. use estimated radio maps to plan paths around communication risk \cite{gordon2026communication}. The proposed experiment instead isolates whether a near-term link-outage forecast changes online task reservation and survey outcomes, compared with reacting after disconnection.

Coverage planning, communication-aware coordination, and prediction are already investigated. Martorell-Torres et al. also propose a coordination approach using acoustic communication quality for a heterogeneous marine team \cite{martorell2024coordination}; its underwater link is not our above-water boat-to-boat radio model. Gong et al. use prediction of path variables in resource-aware USV path following \cite{gong2022resource}; this is not a forecast of a future radio outage. These distinctions prevent an unsupported claim that the field is empty.

For algorithm provenance, Choi et al.'s CBAA/CBBA provides a published auction-and-consensus allocation reference \cite{choi2009auction}. Our planned allocator is a study-defined deterministic greedy rule based on projected completion time. It does not reproduce CBBA's bundle construction or distributed consensus, and its performance guarantees will not be assumed. PythonRobotics provides publicly inspectable grid-sweep coverage code \cite{sakai2018pythonrobotics}. Its route generator is a candidate for checking or adapting the coverage component, not an existing solution to intermittent USV coordination. The experiment compares reservation policies under a common allocator; it does not establish superiority over published task-allocation algorithms.
